\section{Anhang}

\subsection{Massenträgheitsmomente}\label{Massenträgheitsmomente}

La linea rossa tratteggiata è l'asse di rotazione (\textit{Drehachse}).
\smallbreak

\begingroup
\setlength{\extrarowheight}{5pt}
\renewcommand{\arraystretch}{1.25}
\begin{ctabular}{| c | c |}
    \hline
    \textbf{Zylinder / Scheibe}                                                 & \textbf{Hohlzylinder}                                                     \\
    \begin{tikzpicture}[>=Latex, scale=0.9]
  % Körper
  \path[left color=gray!70, right color=gray!25, draw=black, thick]
    (0,0.8) -- (2,0.8) arc(90:-90:0.3 and 0.8) -- (0,-0.8) arc(270:90:0.3 and 0.8) -- cycle;
  \draw[thick, fill=gray!10] (2,0) ellipse (0.3 and 0.8);

  % Drehachse
  \draw[red, dashed, thick] (-0.9,0) -- (3.1,0);

  % Radius
  \draw[<->] (2,0) -- (2,0.8) node[midway, right=1pt] {\footnotesize $r$};
\end{tikzpicture}
             & \begin{tikzpicture}[>=Latex, scale=0.9]
  % Körper
  \path[left color=gray!70, right color=gray!25, draw=black, thick]
    (0,0.8) -- (2,0.8) arc(90:-90:0.3 and 0.8) -- (0,-0.8) arc(270:90:0.3 and 0.8) -- cycle;
  \draw[thick, fill=gray!10] (2,0) ellipse (0.3 and 0.8);
  \draw[thick, fill=gray!60] (2,0) ellipse (0.15 and 0.4);

  % Drehachse
  \draw[red, dashed, thick] (-0.9,0) -- (3.1,0);

  % Radien
  \draw[<->] (2,0) -- (2,0.4) node[midway, right=1pt] {\footnotesize $r_1$};
  \draw[<->] (2.5,0) -- (2.5,0.8) node[midway, right=1pt] {\footnotesize $r_2$};
  \draw[dotted] (2,0.8) -- (2.5,0.8);
\end{tikzpicture}
       \\
    $J = \frac{1}{2} m r^2$                                                     & $J = m \, \frac{r_1^2 + r_2^2}{2}$                                        \\
    \hline
    \textbf{Dünner Stab (Mitte)}                                                & \textbf{Dünner Stab (Ende)}                                               \\
    \begin{tikzpicture}[>=Latex, scale=0.9]
  % Stab
  \path[top color=gray!25, bottom color=gray!70, draw=black, thick, rounded corners=1.5pt]
    (0,-0.12) rectangle (3,0.12);

  % Drehachse
  \draw[red, dashed, thick] (1.5,-0.7) -- (1.5,0.8);

  % Länge
  \draw[<->] (0,0.4) -- (3,0.4) node[midway, fill=white, inner sep=1pt] {\footnotesize $l$};
\end{tikzpicture}
   & \begin{tikzpicture}[>=Latex, scale=0.9]
  % Stab
  \path[top color=gray!25, bottom color=gray!70, draw=black, thick, rounded corners=1.5pt]
    (0,-0.12) rectangle (3,0.12);

  % Drehachse
  \draw[red, dashed, thick] (3,-0.7) -- (3,0.8);

  % Länge
  \draw[<->] (0,0.4) -- (3,0.4) node[midway, fill=white, inner sep=1pt] {\footnotesize $l$};
\end{tikzpicture}
  \\
    $J = \frac{1}{12} m l^2$                                                    & $J = \frac{1}{3} m l^2$                                                   \\
    \hline
    \textbf{Kugel}                                                              & \textbf{Massenpunkt}                                                      \\
    \begin{tikzpicture}[>=Latex, scale=0.9]
  % Kugel
  \shade[ball color=gray!45] (0,0) circle (0.9);
  \draw[thick] (0,0) circle (0.9);

  % Drehachse
  \draw[red, dashed, thick] (0,-1.2) -- (0,1.2);

  % Radius
  \filldraw (0,0) circle (0.03);
  \draw[<->] (0,0) -- (0.9,0) node[midway, above] {\footnotesize $r$};
\end{tikzpicture}
                & \begin{tikzpicture}[>=Latex, scale=0.9]
  % Drehachse
  \draw[red, dashed, thick] (0,-0.8) -- (0,0.8);

  % Masse
  \shade[ball color=gray!50] (1.6,0) circle (0.18);
  \node[below=2pt] at (1.6,-0.18) {\footnotesize $m$};

  % Abstand
  \draw[<->] (0,0) -- (1.42,0) node[midway, above] {\footnotesize $r$};
\end{tikzpicture}
         \\
    $J = \frac{2}{5} m r^2$                                                     & $J = m r^2$                                                               \\
    \hline
\end{ctabular}
\endgroup


\subsection{Analogie Translation \texorpdfstring{$\leftrightarrow$}{<->} Rotation}\label{AnalogieRotation}

\begingroup
\renewcommand{\arraystretch}{1.4}
\begin{ctabular}{| l | c | c | l |}
    \hline
    \textbf{Translation}     & & & \textbf{Rotation}                         \\
    \hline
    Weg                      & $x$                  & $\varphi$                  & Winkel                    \\
    Geschwindigkeit          & $v = \dot{x}$        & $\omega = \dot{\varphi}$   & Winkelgeschwindigkeit     \\
    Beschleunigung           & $a = \ddot{x}$       & $\alpha = \ddot{\varphi}$  & Winkelbeschleunigung      \\
    Masse                    & $m$                  & $J$                        & Trägheitsmoment           \\
    Kraft                    & $F = m\ddot{x}$      & $M = J\ddot{\varphi}$      & Drehmoment                \\
    Federkonstante           & $c$                  & $c^*$                      & Winkelrichtgrösse         \\
    Rückstellkraft           & $F = -c\,x$          & $M = -c^*\varphi$          & Rückstellmoment           \\
    Dämpfung                 & $b\,\dot{x}$         & $b^*\dot{\varphi}$         & Drehdämpfung              \\
    Impuls                   & $p = m v$            & $L = J\omega$              & Drehimpuls                \\
    kin. Energie             & $\tfrac{1}{2}mv^2$   & $\tfrac{1}{2}J\omega^2$    & Rotationsenergie          \\
    pot. Energie             & $\tfrac{1}{2}c\,x^2$ & $\tfrac{1}{2}c^*\varphi^2$ & pot. Energie (Torsion)    \\
    \hline
\end{ctabular}
\endgroup

\medskip

\textbf{Schwingungen:} gleiche DGL, nur Grössen austauschen
\begin{align*}
    m\ddot{x} + b\dot{x} + c\,x &= F_0 \sin(\omega t)
    & \omega_0^2 &= \tfrac{c}{m} & \delta &= \tfrac{b}{2m} \\
    J\ddot{\varphi} + b^*\dot{\varphi} + c^*\varphi &= M_0 \sin(\omega t)
    & \omega_0^2 &= \tfrac{c^*}{J} & \delta &= \tfrac{b^*}{2J}
\end{align*}


\subsubsection{Drehpendel mit Federerregung (Pohl)}\label{PohlFedererregung}

\textrightarrow\ Motor dreht über Exzenter und Schubstange den \textbf{Hebel} (= äusseres Federende) um $\varphi_e(t) = \varphi_{e,0} \sin(\omega t)$

\centerline{\begin{tikzpicture}[font=\footnotesize, >={Stealth}, scale=0.9]
  % --- Ruhelage ---
  \draw[gray, densely dashed] (0,0) -- (0,2.6);
  \node[gray, above] at (0,2.6) {Ruhelage};
  % --- Rad (J) mit Speichen, um phi = 15 Grad gedreht ---
  \draw[blue!45, line width=3pt] (0,0) circle (1.3);
  \foreach \a in {15,75,135,195,255,315}
    \draw[blue!25, thick] (0,0) -- (\a:1.2);
  \draw[blue, line width=2pt] (75:1.05) -- (75:1.6);
  % --- Spiralfeder c*: innen am Rad, aussen am Hebel ---
  \draw[gray!80!black, thick, domain=0:1, samples=160, smooth, variable=\x]
    plot ({(0.15+0.55*\x)*cos(50-900*(1-\x))}, {(0.15+0.55*\x)*sin(50-900*(1-\x))});
  % --- Hebel, um phi_e = 40 Grad gedreht ---
  \draw[orange, line width=2.5pt, line cap=round] (230:2.0) -- (50:2.15);
  \fill[orange] (50:0.7) circle (1.8pt);
  \fill[gray!70!black] (0,0) circle (2pt);
  % --- Winkel ab Ruhelage ---
  \draw[blue, thick, ->] (90:1.85) arc (90:75:1.85);
  \node[blue] at (82:2.12) {$\varphi$};
  \draw[orange, thick, ->] (90:2.45) arc (90:50:2.45);
  \node[orange] at (63:2.75) {$\varphi_e$};
  % --- Federmoment auf das Rad ---
  \draw[violet, very thick, ->] (72:1.48) arc (72:54:1.48);
  \node[violet, anchor=west] at (25:1.55) {$M = c^*(\varphi_e - \varphi)$};
  % --- Schubstange und Motor mit Exzenter ---
  \draw[gray!70!black, line width=1.5pt] (230:2.0) -- (2.3,-1.532);
  \draw[thick, fill=gray!15] (2.7,-1.532) circle (0.45);
  \fill (2.3,-1.532) circle (1.8pt);
  \fill[orange] (230:2.0) circle (1.8pt);
  \draw[->] (3.27,-1.2) arc (35:-60:0.6);
  \node at (3.7,-1.75) {$\omega$};
  % --- Beschriftungen ---
  \node[blue, anchor=west] at (1.2,-0.7) {Rad $J$};
  \node[gray!70!black, anchor=east] at (-0.8,0.35) {Feder $c^*$};
  \node[orange, anchor=east] at (-1.35,-1.55) {Hebel};
  \node at (0.6,-1.85) {Schubstange};
  \node at (2.7,-2.2) {Motor};
\end{tikzpicture}
}

\smallskip

Die Feder wirkt nur über ihre \textbf{Verdrillung} $\varphi_e - \varphi$: \quad $M = -c^*(\varphi - \varphi_e)$

$$ \boxed{ J\ddot{\varphi} + b^*\dot{\varphi} + c^*\varphi = c^* \varphi_{e,0} \sin(\omega t) } \qquad \rightarrow \quad M_0 = c^* \varphi_{e,0} $$

$$ \boxed{ \frac{A}{\varphi_{e,0}} = V } \qquad \text{Resonanz:} \quad \frac{A}{\varphi_{e,0}} = V_r \qquad (V,\ V_r \text{ siehe Abschnitt \ref{Vergrösserungsfunktion / Phasenverschiebung}}) $$

\textrightarrow\ $J$, $c^*$ und $M_0$ kürzen sich weg: es braucht nur $D$ (aus $T$ und Abklingen)

\begin{ctabular}{c l c}
    $\varphi$         & Winkel des Rades                       & $[\varphi] = \rad$          \\
    $\varphi_e$       & Winkel des Hebels (Erreger)            & $[\varphi_e] = \rad$        \\
    $\varphi_{e,0}$   & Amplitude des Erregers                 & $[\varphi_{e,0}] = \rad$    \\
    $A$               & Amplitude des Rades                    & $[A] = \rad$                \\
    $c^*$             & Winkelrichtgrösse der Spiralfeder      & $[c^*] = \frac{\newton\meter}{\rad}$
\end{ctabular}


\subsubsection{DGL aufstellen: le regole dei segni}\label{DGLVorzeichen}

\centerline{\begin{tikzpicture}[font=\scriptsize, >={Stealth}]
  % Feder horizontal von #1 nach #2 auf Höhe #3
  \def\hfeder#1#2#3{\draw[thick] (#1,#3) -- (#1+0.12,#3)
      plot[domain=0:1, samples=21, variable=\t] ({#1+0.12+(#2-#1-0.24)*\t}, {#3+0.12*asin(sin(\t*1800))/90})
      -- (#2,#3);}
  % Wand
  \fill[gray!35] (-0.2,-0.6) rectangle (0,0.6);
  % k1
  \hfeder{0}{1.4}{0}
  \node[above] at (0.7,0.12) {$k_1$};
  % m1
  \draw[thick, fill=blue!15] (1.4,-0.4) rectangle (2.1,0.4);
  \node at (1.75,0) {$m_1$};
  % k zwischen m1 und m2
  \hfeder{2.1}{3.9}{0.22}
  \node[above] at (3.0,0.3) {$k$};
  % Dämpfer b
  \draw[thick] (2.1,-0.22) -- (2.85,-0.22);
  \draw[thick] (2.7,-0.38) -- (3.15,-0.38) -- (3.15,-0.06) -- (2.7,-0.06);
  \draw[thick] (2.85,-0.34) -- (2.85,-0.10);
  \draw[thick] (3.15,-0.22) -- (3.9,-0.22);
  \node[below] at (3.0,-0.38) {$b$};
  % m2
  \draw[thick, fill=blue!15] (3.9,-0.4) rectangle (4.6,0.4);
  \node at (4.25,0) {$m_2$};
  % externe Kraft
  \draw[->, very thick, violet] (4.6,0) -- (5.4,0) node[right] {$F(t)$};
  % Kräfte auf m1 (m1 um +x1 verschoben gedacht)
  \draw[->, very thick, blue] (1.75,0.6) -- (1.05,0.6);
  \node[blue, anchor=south east] at (1.4,0.62) {$-k_1 x_1$};
  \draw[->, very thick, orange] (1.75,0.6) -- (2.45,0.6);
  \node[orange, anchor=south west] at (2.0,0.62) {$+k\,(x_2-x_1)$};
  \draw[->, very thick, orange] (1.75,-0.6) -- (2.45,-0.6);
  \node[orange, anchor=north west] at (1.55,-0.68) {$+b\,(\dot x_2-\dot x_1)$};
  % Achse
  \draw[->, gray!60!black, thick] (0,-1.25) -- (5.6,-1.25) node[right] {$+x$};
  \node[gray!60!black, anchor=west] at (0,-1.45) {\textbf{1.} Achse zuerst};
\end{tikzpicture}
}

\smallskip
{\small Idea: \textbf{sposta la massa di $+x$} (e dalle velocità $+\dot{x}$) e chiediti per ogni forza: \emph{la riporta indietro o la spinge avanti?}}

\renewcommand{\arraystretch}{1.35}
\begin{tabular}{@{}>{\raggedright\arraybackslash}p{0.25\columnwidth}>{\raggedright\arraybackslash}p{0.71\columnwidth}@{}}
    \textbf{1. Asse} & disegna la freccia di $x$ (o il verso di $\varphi$): forza nel verso della freccia $= +$, contro $= -$ \\
    \textbf{2. Molla verso il muro} & riporta indietro $\Rightarrow$ \cbl{$-k\,x$} \\
    \textbf{3. Molla fra $i$ e $j$} & allungata tira verso l'altra estremità. \;Su $i$: \cor{$+k\,(x_j - x_i)$}, su $j$: \cor{$+k\,(x_i - x_j)$} \\
    \textbf{4. Smorzatore fra $i$ e $j$} & come la molla, con le velocità. \;Su $i$: \cor{$+b\,(\dot{x}_j - \dot{x}_i)$} \;(verso il muro: $-b\,\dot{x}$) \\
    \textbf{5. Forza costante} & $mg$, \ldots: $+$ se nel verso dell'asse. Sposta solo $x_0$, non cambia $\omega$ \\
    \textbf{6. Eccitazione $u(t)$} & il supporto mobile è la ``massa $j$'' delle regole 3 e 4: \;$k\,(u - x)$, \;$b\,(\dot{u} - \dot{x})$ \\
    \textbf{7. Rotazione} & stesse regole con $\varphi$, $M$, $J$: \;$M = F \cdot$ braccio, $+$ se fa ruotare nel verso di $\varphi$
\end{tabular}
\renewcommand{\arraystretch}{1}

\smallskip
\noindent\fbox{\parbox{0.96\columnwidth}{%
\textbf{Controlli finali} \quad porta tutto a sinistra: $\;m_i\ddot{x}_i + \ldots + (\ldots)\,x_i - (\ldots)\,x_j = F_i$
\begin{itemize}
    \item[\textbf{a}] coefficienti di $\ddot{x}_i$, $\dot{x}_i$, $x_i$ \textbf{tutti positivi} \;(altrimenti: segno sbagliato)
    \item[\textbf{b}] coefficiente di $x_i$ $=$ \textbf{somma delle molle attaccate} a $i$ \;(es. $k_1 + k$)
    \item[\textbf{c}] \textbf{accoppiamento simmetrico}: $-k\,x_j$ nell'eq.\ di $i$ $\Leftrightarrow$ $-k\,x_i$ nell'eq.\ di $j$
\end{itemize}}}

\smallskip
\textbf{Esempio (figura):}
\begin{align*}
    m_1 \ddot{x}_1 + b\,\dot{x}_1 + (k_1 + k)\,x_1 - b\,\dot{x}_2 - k\,x_2 &= 0 \\
    m_2 \ddot{x}_2 + b\,\dot{x}_2 + k\,x_2 - b\,\dot{x}_1 - k\,x_1 &= F(t)
\end{align*}


\subsection{Trigonometrie}

\begingroup
\renewcommand{\arraystretch}{2}
\setlength{\tabcolsep}{0mm}
\scalebox{0.33}{
\Huge
\begin{tabularx}{3\columnwidth}{rp{1mm}V{3}*{17}{C}}
    \rowcolor{subsectioncolor!30}$\bm{\alpha}$\phantom{${}^{\circ}$} && $0$ & $\dfrac{\pi}{6\mathstrut}$ & $\dfrac{\pi}{4}$ & $\dfrac{\pi}{3}$ & $\dfrac{\pi}{2}$ & $\dfrac{2\pi}{3}$ & $\dfrac{3\pi}{4}$ & $\dfrac{5\pi}{6}$ & $\pi$ & $\dfrac{7\pi}{6}$ & $\dfrac{5\pi}{4}$ & $\dfrac{4\pi}{3}$ & $\dfrac{3\pi}{2}$ & $\dfrac{5\pi}{3}$ & $\dfrac{7\pi}{4}$ & $\dfrac{11\pi}{6}$ & $2\pi$ \\\hline
    \rowcolor{subsectioncolor!30}$\bm{\alpha^\circ}$ && \phantom{${}^\circ$}$0^\circ$ & $30^\circ$ & $45^\circ$ & $60^\circ$ & $90^\circ$ & $120^\circ$ & $135^\circ$ & $150^\circ$ & $180^\circ$ & $210^\circ$ & $225^\circ$ & $240^\circ$ & $270^\circ$ & $300^\circ$ & $315^\circ$ & $330^\circ$ & $360^\circ$ \\\hline
    $\bm{\sin(\alpha)}$ && $0$ & $\dfrac{1}{2\mathstrut}$ & $\dfrac{\sqrt{2}}{2}$ & $\dfrac{\sqrt{3}}{2}$ & $1$ & $\dfrac{\sqrt{3}}{2}$ & $\dfrac{\sqrt{2}}{2}$ & $\dfrac{1}{2}$ & $0$ & $-\dfrac{1}{2}$ & $-\dfrac{\sqrt{2}}{2}$ & $-\dfrac{\sqrt{3}}{2}$ & $-1$ & $-\dfrac{\sqrt{3}}{2}$ & $-\dfrac{\sqrt{2}}{2}$ & $-\dfrac{1}{2}$ & $0$ \\\hline
    $\bm{\cos(\alpha)}$ && $1$ & $\dfrac{\sqrt{3}}{2\mathstrut}$ & $\dfrac{\sqrt{2}}{2}$ & $\dfrac{1}{2}$ & $0$ & $-\dfrac{1}{2}$ & $-\dfrac{\sqrt{2}}{2}$ & $-\dfrac{\sqrt{3}}{2}$ & $-1$ & $-\dfrac{\sqrt{3}}{2}$ & $-\dfrac{\sqrt{2}}{2}$ & $-\dfrac{1}{2}$ & $0$ & $\dfrac{1}{2}$ & $\dfrac{\sqrt{2}}{2}$ & $\dfrac{\sqrt{3}}{2}$ & $1$ \\\hline
    $\bm{\tan(\alpha)}$ && $0$ & $\dfrac{\sqrt{3}}{3\mathstrut}$ & $1$ & $\sqrt{3}$ & $\pm\infty$ & $-\sqrt{3}$ & $-1$ & $-\dfrac{\sqrt{3}}{3}$ & $0$ & $\dfrac{\sqrt{3}}{3}$ & $1$ & $\sqrt{3}$ & $\pm\infty$ & $-\sqrt{3}$ & $-1$ & $-\dfrac{\sqrt{3}}{3}$ & $0$ \\\hline
    $\bm{\cot(\alpha)}$ && $\pm\infty$ & $\sqrt{3}$ & $1$ & $\dfrac{\sqrt{3}}{3\mathstrut}$ & $0$ & $-\dfrac{\sqrt{3}}{3}$ & $-1$ & $-\sqrt{3}$ & $\pm\infty$ & $\sqrt{3}$ & $1$ & $\dfrac{\sqrt{3}}{3}$ & $0$ & $-\dfrac{\sqrt{3}}{3}$ & $-1$ & $-\sqrt{3}$ & $\pm\infty$ \\
\end{tabularx}}
\endgroup

\[\sin(\alpha) = \frac{\rm C. Opposto}{\rm Ipotenusa} \qquad 
\cos(\alpha) = \frac{\rm C. Adiacente}{\rm Ipotenusa} \qquad 
\tan(\alpha) = \frac{\rm C. Opposto}{\rm C. adiacente}\]


\subsubsection{Beziehungen zwischen $\bm{\sin(x)}$ und $\bm{\cos(x)$}}

\begin{tabular}{lll}
    $\sin(-a) = -\sin(a)$   & $\sin(\pi - a) = \sin(a)$     & $\sin(\pi + a) = -\sin(a)$ \\	
    $\cos(-a) = \cos(a)$    & $\cos(\pi - a) = - \cos(a)$   & $\cos(\pi + a) = - \cos(a)$ \\
    \\
\end{tabular}

\begin{tabular}{ll}
    $\sin(\frac{\pi}{2} - a) = \sin(\frac{\pi}{2} + a) = \cos(a)$ & $\cos(\frac{\pi}{2} - a) = - \cos(\frac{\pi}{2} + a) = \sin(a)$		 
\end{tabular}


\subsubsection{Additionstheoreme}

\begin{tabular}{ll}
    $\sin(a \pm b) = \sin(a) \cdot \cos(b) \pm \cos(a) \cdot \sin(b)$   & $\tan(a \pm b) = \frac{\tan(a) \pm \tan(b)}{1 \mp \tan(a) \cdot \tan(b)}$\\
    $\cos(a \pm b) = \cos(a) \cdot \cos(b) \mp \sin(a) \cdot \sin(b)$   &
\end{tabular}


\subsubsection{Summen und Differenzen}

\renewcommand{\arraystretch}{1.5}
\begin{tabular}{ll@{}}
    $\sin(a) + \sin(b) = 2 \cdot \sin \big( \frac{a+b}{2} \big) \cdot \cos \big( \frac{a-b}{2} \big)$ &
    $\cos(a) + \cos(b) = 2 \cdot \cos \big( \frac{a+b}{2} \big) \cdot \cos \big( \frac{a-b}{2} \big)$ \\
    $\sin(a) - \sin(b) = 2 \cdot \sin \big( \frac{a-b}{2} \big) \cdot \cos \big( \frac{a+b}{2} \big)$ &
    $\cos(a) - \cos(b) = -2 \cdot \sin \big( \frac{a+b}{2} \big) \cdot \sin \big( \frac{a-b}{2} \big)$ \\
    $\tan(a) \pm \tan(b) = \frac{\sin(a \pm b)}{\cos(a) \cdot \cos(b)}$ 
\end{tabular}
\renewcommand{\arraystretch}{1}


\subsubsection{Produkte}

$\sin(a) \cdot \sin(b) = \frac{1}{2} \big( \cos(a-b) - \cos(a+b) \big) $ \\
$\cos(a) \cdot \cos(b) = \frac{1}{2} \big( \cos(a-b) + \cos(a+b) \big) $ \\	
$\sin(a) \cdot \cos(b) = \frac{1}{2} \big( \sin(a-b) + \sin(a+b) \big) $ 
    
    
\subsubsection{Winkelvielfache und Halbwinkel}

\renewcommand{\arraystretch}{1.3}
\begin{tabular}{ll}
    $\sin(2a) = 2 \, \sin(a) \cdot \cos(a)$                         & $\cos(2a) = \cos^2(a) - \sin^2(a)$                                \\
    $\sin(3a) = 3 \, \sin(a) -4 \, \sin^3(a)$                       & $\cos(3a) = 4 \, \cos^3(a) -3 \, \cos(a)$                         \\ 
    $\sin \big( \frac{a}{2} \big) = \sqrt{\frac{1}{2}(1-\cos(a))} $ & $\cos \big( \frac{a}{2} \big) = \sqrt{\frac{1}{2}(1+\cos(a))} $ 
\end{tabular}


\subsubsection{Potenzen}

\renewcommand{\arraystretch}{1.5}
\begin{tabular}{ll}
    $\sin^2(a) = \frac{1^{\mathstrut}}{2}(1-\cos(2a))$               & $\cos^2(a) = \frac{1}{2}(1+\cos(2a))$                \\
    $\sin^3(a) = \frac{1^{\mathstrut}}{4}(3\, \sin(a) - \sin(3a))$   & $\cos^3(a) = \frac{1}{4}(\cos(3a) + 3 \, \cos(a))$
\end{tabular}
\renewcommand{\arraystretch}{1}

\subsection{Fehlerfortpflanzung}

Data una grandezza calcolata da più variabili misurate, es.\ $f = f(c,\lambda) = \dfrac{c}{\lambda}$, sia $c$ che $\lambda$ portano un errore ($\Delta c$, $\Delta\lambda$) che si propaga sul risultato $\Delta f$.

\textbf{Contributi singoli:}
\[
\frac{\partial f}{\partial \lambda} = -\frac{c}{\lambda^2}
\quad\Rightarrow\quad
\Delta f_\lambda = \frac{\partial f}{\partial \lambda}\,\Delta\lambda = -\frac{c}{\lambda^2}\,\Delta\lambda
\]
\[
\frac{\partial f}{\partial c} = \frac{1}{\lambda}
\quad\Rightarrow\quad
\Delta f_c = \frac{\partial f}{\partial c}\,\Delta c = \frac{1}{\lambda}\,\Delta c
\]

In generale, per $f = f(x_1, x_2, \dots, x_n)$:
\[
\Delta f = \sqrt{\sum_{i=1}^n \left(\frac{\partial f}{\partial x_i}\right)^2 \Delta x_i^2}
\]